\slajdid{10-s036}
\begin{frame}[label=10-s036]
\gusttekst\setlength{\abovedisplayskip}{3pt}\setlength{\belowdisplayskip}{3pt}\setlength{\abovedisplayshortskip}{2pt}\setlength{\belowdisplayshortskip}{2pt}Situacija kada nema otpornika koji će definisati napon logičke jedinice – stanje visoke impedanse
\par\medskip
Problem koji se pojavljuje jeste pitanje: Šta se nalazi na liniji kada su svi izlazni tranzistori logičkih kola 1, 2 i 3 zakočeni. \underline{Odgovor ništa ne postoji u elektronici.} Kako to shvata naredno logičko kolo 4. Kao prvo, jednačina Bulove algebre više ne važi pošto ne postoji stanje logičke jedinice na liniji. U digitalnoj elektronici to stanje se naziva stanjem visoke impedanse.
\begin{columns}[c,onlytextwidth]\column{.49\textwidth}\centering\begin{tikzpicture}[circuit logic US,x=.9cm,y=.75cm,font=\sitneoznake,>=Latex]\ctikzset{bipoles/length=.55cm}\node[not gate US,draw,minimum width=9mm,minimum height=7mm](a)at(0,1.5){1};\node[nand gate US,draw,logic gate inputs=nn,minimum width=9mm,minimum height=9mm](b)at(0,0){2};\node[or gate US,draw,logic gate inputs=nn,minimum width=9mm,minimum height=9mm](c)at(0,-1.5){3};\draw(a.input)--++(-.6,0)node[left]{A};\draw(b.input 1)--++(-.6,0)node[left]{B};\draw(c.input 2)--++(-.6,0)node[left]{D};\draw(b.input 2)--++(-.25,0)coordinate(bc);\draw(c.input 1)--++(-.25,0)coordinate(cc);\draw(bc)--(cc);\draw($(bc)!.5!(cc)$)--++(-.35,0)node[left]{C};\draw(a.output)--(1.7,1.5);\node[above]at(.9,1.5){oc};\draw(b.output)--(1.7,0);\node[above]at(.9,0){oc};\draw(c.output)--(1.7,-1.5);\node[above]at(.9,-1.5){oc};\draw(1.7,1.5)--(1.7,-1.5);\draw(1.7,-.6)--(2.6,-.6);\node[not gate US,draw,minimum width=9mm,minimum height=7mm](d)at(3.8,-.6){4};\draw(2.6,-.6)--(d.input);\draw(d.output)--++(.35,0);\node at(2.6,.2){$(\bar A)(\overline{BC})(CD)$};\draw(1.9,-.15)--(3.3,.55);\draw(1.9,.55)--(3.3,-.15);\end{tikzpicture}
\column{.49\textwidth}\centering\begin{tikzpicture}[circuit logic US,x=.9cm,y=.75cm,font=\sitneoznake,>=Latex]\ctikzset{bipoles/length=.55cm}\node[not gate US,draw,minimum width=9mm,minimum height=7mm](a)at(0,1.5){1};\node[nand gate US,draw,logic gate inputs=nn,minimum width=9mm,minimum height=9mm](b)at(0,0){2};\node[or gate US,draw,logic gate inputs=nn,minimum width=9mm,minimum height=9mm](c)at(0,-1.5){3};\draw(a.input)--++(-.6,0)node[left]{A};\draw(b.input 1)--++(-.6,0)node[left]{B};\draw(c.input 2)--++(-.6,0)node[left]{D};\draw(b.input 2)--++(-.25,0)coordinate(bc);\draw(c.input 1)--++(-.25,0)coordinate(cc);\draw(bc)--(cc);\draw($(bc)!.5!(cc)$)--++(-.35,0)node[left]{C};\draw(a.output)--(1.7,1.5);\node[above]at(.9,1.5){oc};\draw(b.output)--(1.7,0);\node[above]at(.9,0){oc};\draw(c.output)--(1.7,-1.5);\node[above]at(.9,-1.5){oc};\draw(1.7,1.5)--(1.7,-1.5);\draw(1.7,-.6)--(2.6,-.6);\node[not gate US,draw,minimum width=9mm,minimum height=7mm](d)at(3.8,-.6){4};\draw(2.6,-.6)--(d.input);\draw(d.output)--++(.35,0);\draw(3.1,-1.9)--(3.1,-1.4)--(2.8,-1.4)--(2.8,-1.6)--(2.5,-1.25)--(2.8,-.9)--(2.8,-1.1)--(3.1,-1.1);\node[right]at(3.1,-1.8){$Z\to\infty$};\end{tikzpicture}
\end{columns}
\par\medskip Impedansa koju vidi ulaz u logičko kolo 4 prema liniji spajanja, odnosno prema izlazima logičkih kola 1, 2 i 3 kada su svi izlazni tranzistori u njima zakočeni je beskonačna. Nema putanje prema napajanju ili masi.
\end{frame}
